\section{模型评价与推广}
\subsection{模型优点}
第一，任务瞬时容量与小时电量使用不同但一致的重叠口径，既不会用平均负荷掩盖短时 GPU 峰值，也不会把分钟级任务按整小时计费。第二，预测严格按训练、验证、测试顺序运行，多步测试不使用未来真值；季节朴素基线与 Huber 网格的比较保留了可解释参照。第三，Q2 的增量搜索与全任务审计分离，使 50000 次候选评价仍能保留五次独立全量约束确认。第四，Q3 强制 no-action 入选并对轨迹去重，避免只展示储能方案而缺少基线。第五，Q4 每个情景都动态重建邻域，不把固定任务候选误当成联合优化。

结果交付保持了“数值—数据表—约束审计—图形”的可追踪链。五方案、七策略、17 情景均有逐行 CSV；两次隔离运行还验证了除自然耗时外的字节级一致性。这些证据支持方案属于可复算的场景可行候选。

\subsection{模型不足}
最主要限制是有限搜索深度。Q2 每个方向只检查 10000 个单任务移动，Q4 每轮只检查 600 个候选并最多运行 8 轮；即使全量硬约束审计通过，也无法排除未访问邻域中存在更优解。Q3 的 21 状态 SOC 离散会引入网格近似，候选集内非支配不等同于连续空间 Pareto 前沿。

储能模型未包含退化成本、循环寿命、自放电和备用容量。平衡策略的等效循环合计超过 1500，表明账面收益可能依赖过高频率的充放电。成本又按购电支出减售电收入计算，且售电规模远大于购电量；负成本应解释为题面边界下的净现金流，而非系统能耗或全生命周期成本为负。

通信侧只使用静态单向时延，未考虑带宽、拥塞、迁移数据量与传输能耗。服务质量包含等待代理，平衡方案最大等待达到 485 小时，说明平均或归一化服务指标不能代替任务级最长等待检查。碳约束 20\% 与 30\% 未达到也只反映当前启发式，不是数学不可行性证明。

最后，当前 benchmark 没有独立 Oracle，结果只能达到 scenario-feasible。双运行和约束残差能证明可重复与内部一致，却不能证明外部真实最优。

\subsection{推广方向}
若扩大计算预算，优先升级 Q2/Q4 的邻域而非重写数据管线：保留现有可行构造作为热启动，引入大邻域搜索、混合整数规划上界或分解协调，并继续使用同一全量审计。Q3 可在现有状态转移中加入每 MWh 吞吐退化成本和年度循环上限，再进行网格加密对照。

面向在线运行，可采用滚动时域：每小时更新任务、新能源和价格预测，冻结已开工任务并传递绝对 SOC；这一结构与模型预测控制的状态滚动思想一致\cite{rawlings2017model}。多目标比较则可在统一尺度下采用 $\varepsilon$-约束法生成更充分的非支配候选\cite{miettinen1999nonlinear}，同时保留“候选内非支配”和“已证 Pareto”之间的表述边界。
